% Build in this directory: pdflatex -interaction=nonstopmode -halt-on-error class2_handout.tex
\documentclass[a4paper,10pt,pdflatex,ja=standard,jafont=ipaex-type1,
  magstyle=real,scale=1]{bxjsarticle}
\usepackage{lmodern}
\usepackage{amsmath,amssymb,mathtools}
\usepackage{geometry}
\geometry{reset,left=22mm,right=22mm,top=20mm,bottom=17mm,
  headheight=13pt,headsep=6mm,footskip=10mm}
\usepackage{xcolor}
\usepackage{enumitem}
\usepackage{array,booktabs,tabularx}
\usepackage{fancyhdr}

\definecolor{ink}{HTML}{193A5C}
\definecolor{accent}{HTML}{237D82}
\definecolor{soft}{HTML}{EAF3F3}
\definecolor{muted}{HTML}{52616B}
\newcommand{\coursename}{解析学 II}
\newcommand{\handoutnumber}{2}
\pagestyle{fancy}
\fancyhf{}
\fancyhead[L]{\small\color{ink}\coursename{}・授業要約と演習}
\fancyhead[R]{\small\color{ink}第\handoutnumber 回}
\fancyfoot[C]{\small\color{ink}\thepage/2}
\renewcommand{\headrulewidth}{0.3pt}
\setlength{\parindent}{0pt}
\setlength{\parskip}{2pt}
\setlength{\abovedisplayskip}{5pt}
\setlength{\belowdisplayskip}{5pt}
\setlength{\abovedisplayshortskip}{3pt}
\setlength{\belowdisplayshortskip}{4pt}
\setlist[itemize]{leftmargin=1.8em,itemsep=1pt,topsep=2pt,parsep=0pt}

\newcommand{\handouttitle}[1]{%
  \begin{center}
    {\Large\color{ink}#1}\par\vspace{3mm}
    {\color{accent}\rule{0.88\linewidth}{0.7pt}}
  \end{center}\vspace{-2mm}}
\newcommand{\handout}[2]{%
  \renewcommand{\handoutnumber}{#1}%
  \handouttitle{第#1回\quad #2}}
\newcommand{\topic}[1]{\par\vspace{5pt}%
  {\large\sffamily\mdseries\color{accent}#1}\par\vspace{2pt}}
\newcommand{\keybox}[1]{\par\vspace{3pt}\begingroup
  \setlength{\fboxsep}{4pt}%
  \colorbox{soft}{\parbox{\dimexpr\linewidth-2\fboxsep\relax}{\small #1}}%
  \endgroup\par\vspace{2pt}}
\newenvironment{problems}
  {\begin{enumerate}[leftmargin=2em,itemsep=3pt,topsep=3pt,parsep=0pt,
    font=\color{accent}]}
  {\end{enumerate}}


\newcommand{\examplelabel}[1]{{\sffamily\mdseries\color{ink}#1}}
\newcommand{\answerroom}{\par\vspace{16mm}}

\begin{document}
\fontsize{10pt}{13.5pt}\selectfont
\setlength{\abovedisplayskip}{4pt}
\setlength{\belowdisplayskip}{4pt}
\handout{2}{有理関数の積分}
\topic{有理関数と多項式除法}
実係数多項式 $P,Q$（$Q\not\equiv0$）に対し、$P/Q$ を有理関数という。
積分は $Q(x)\ne0$ となる区間で考える。まず多項式除法により
$$
 \frac{P(x)}{Q(x)}=H(x)+\frac{R(x)}{Q(x)},\qquad \deg R<\deg Q
$$
と書き、多項式部分 $H$ と真分数部分を別々に積分する。

\topic{一次因子・二次因子の基本公式}
$n\ge1$ を整数とすると、$x\ne\alpha$ の区間で
$$
 \int\frac{dx}{(x-\alpha)^n}=
 \begin{cases}
 \log|x-\alpha|+C,&n=1,\\[2pt]
 -\dfrac{1}{(n-1)(x-\alpha)^{n-1}}+C,&n\ge2.
 \end{cases}
$$
$q(x)=x^2+ax+b$、$D=a^2-4b$ とおく。$D>0$ のとき、異なる実根を
$\alpha,\beta$ として $q(x)=(x-\alpha)(x-\beta)$ と書けば
$$
 \int\frac{dx}{q(x)}=
 \begin{cases}
 \dfrac{1}{\alpha-\beta}\log\left|\dfrac{x-\alpha}{x-\beta}\right|+C,&D>0,\\[5pt]
 -\dfrac{1}{x+a/2}+C,&D=0,\\[5pt]
 \dfrac{2}{\sqrt{4b-a^2}}\arctan\dfrac{2x+a}{\sqrt{4b-a^2}}+C,&D<0.
 \end{cases}
$$
二次式に一次式の分子があるときは、$cx+d=\tfrac c2 q'(x)+(d-\tfrac{ac}{2})$
と分け、$q'/q$ の項を対数で積分する。

\topic{部分分数分解}
$Q$ を実数上で一次因子と既約二次因子に分解すると
$$
 \frac{P}{Q}=H+
 \sum_i\sum_{k=1}^{m_i}\frac{A_{ik}}{(x-\alpha_i)^k}
 +\sum_j\sum_{k=1}^{n_j}\frac{B_{jk}x+C_{jk}}{(x^2+a_jx+b_j)^k},
 \qquad a_j^2-4b_j<0.
$$
分母を払って多項式の係数を比較する。重複因子では、1乗から最大のべきまで
すべての項を用意する。

\topic{二次因子の高いべき}
積分定数を最後にまとめ、$I_1(x)=\arctan x$ と選ぶと
$$
 I_n(x)=\frac{x}{2(n-1)(1+x^2)^{n-1}}
       +\frac{2n-3}{2(n-1)}I_{n-1}(x)\quad(n\ge2)
$$
は $(1+x^2)^{-n}$ の原始関数になる。平方完成と一次置換により、
一般の既約二次因子にも使える。

\examplelabel{例：多項式除法から積分する。}
$$
 \frac{x^2}{x^2-1}=1+\frac12\left(\frac1{x-1}-\frac1{x+1}\right),\qquad
 \int\frac{x^2}{x^2-1}\,dx=x+\frac12\log\left|\frac{x-1}{x+1}\right|+C.
$$
\keybox{分母の零点を除いた各区間で原始関数を求める。定積分では、区間内に特異点がないかも確認する。}

\newpage
\handouttitle{第2回\quad 演習問題}
計算過程を記し、必要な条件・積分領域・向きを明示すること。
\begin{problems}
\item $\displaystyle\int\frac{2x+3}{x^2+x-2}\,dx$ を部分分数分解で求めよ。原始関数を考える区間も述べよ。\answerroom
\item $\displaystyle\int\frac{x^3}{x^2+1}\,dx$ を求めよ。まず多項式除法を行うこと。\answerroom
\item $\displaystyle\int\frac{x+1}{x^2+2x+5}\,dx$ を求め、微分して確認せよ。\answerroom
\item 帰納公式を用いて $\displaystyle\int\frac{dx}{(1+x^2)^2}$ を求めよ。\answerroom
\item $\displaystyle\frac{x+1}{x(x-1)^2}=\frac A{x}+\frac B{x-1}+\frac C{(x-1)^2}$ の係数を決め、原始関数を求めよ。\answerroom
\item $\displaystyle\int_0^1\frac{dx}{x^2+2x+2}$ を平方完成と置換積分で求めよ。\answerroom
\end{problems}
\end{document}
